% 本文件由 LaTeX 排版 Agent 根据有效 translation.json 自动生成。
% author=Alexander of Lycopolis; work_id=0618; title=论摩尼教徒

\chapter{论摩尼教徒}

% structure: p0001 source=h1 target=omitted reason=page-title-already-rendered-by-parent

% structure: p0002 source=h2 target=section reason=relative-heading-rank; base=h2

\section{第一章 基督宗教哲学的高超之处；基督徒中异端的起源}

基督徒的哲学被称为简朴。但它非常注重品行的塑造，以隐晦的方式暗示有关天主的更确实的真理；其中最主要的，对于这些事上任何热诚严肃的志向而言，所有人在设定一个极其高贵且极其古老的有效原因，作为一切存在之物的创始者时，都已领受了。因为基督徒将更费力且困难的课题留给伦理学的学生，例如什么是德性，道德的与智性的；以及留给那些花时间构建关于道德、激情与情感的假说，却不指明任何借以获得每种德性的要素，并且仿佛随意堆砌一些不太精细的规诫的人——普通民众听到这些，正如我们由经验所知，便在谦逊上大有长进，一种虔诚的品质被印刻在他们的品行上，使由这些习俗所形成的道德倾向变得活泼，并逐步引导他们向往可敬与良善之事。\par

但这门哲学被后来者的众多问题所分割，便产生了许多派别，正如那些致力于辩证法的人一样，有些人比另一些人更熟练，并且可以说更精于处理微妙细致的问题；以至于现在他们以宗派和异端的创立者与鼻祖的身份出现。由于他们，品行的塑造受到阻碍并变得隐晦；因为那些想要成为宗派首领的人，并未达到话语的确实真理，而普通民众则更被激动起纷争与对抗。而且，对于所争议之事，没有一条规则或法律可以获得解决，相反，如同在其他事务上一样，这种好胜的争竞越演越烈，没有什么不受到它的损害和伤害。\par

% structure: p0005 source=h2 target=section reason=relative-heading-rank; base=h2

\section{第二章 \Person{摩尼}，或称\Person{玛涅斯}的年代；他的首批门徒；二元教义；摩尼教的物质}

因此在这些事上也是这样，人人都在见解的新颖上竭力显示自己为先、为优，竟将这门简朴的哲学几乎化为虚无。他们称这样的人为\Person{摩尼}，是一名波斯人，我在那种教义上的导师是一个名叫\Person{帕普斯}的人，在他之后是\Person{托马斯}，并有其他一些人跟随他们。据说这人生活在\Person{瓦莱里安}皇帝的时代，并效力于波斯王\Person{沙普尔}，因在某事上触怒了他而被处死。关于他的品格和名声，有些这样的传闻从他亲近的人那里传至我处。他设立了两种\OriginalTerm{本原}{principles}：天主与\OriginalTerm{物质}{matter}。他称天主为善，并断言物质为恶。但天主在善上的超越，胜过物质在恶上的程度。但他所说的物质，既不是\Person{柏拉图}所称的——当它接受了性质与形状后便成为万物，因此\Person{柏拉图}称之为全受者，万物的母亲和养育者——也不是\Person{亚里士多德}称之为元素的，与形式和缺失相关的东西，而是除此之外的某种东西。因为那在个别事物中不复合的运动，他就称之为物质。在天主那边排列着如同侍女般的能力，全都是善的；同样，在物质那边也排列着其他能力，全都是恶的。再者，光明、光辉和在上者，这些都与天主同在；而幽暗、黑暗和在下者，则与物质同在。天主也有欲望，但全都是善的；物质同样有欲望，但全都是恶的。\par

% structure: p0007 source=h2 target=section reason=relative-heading-rank; base=h2

\section{第三章 \Person{摩尼}关于\OriginalTerm{物质}{matter}的幻想}

有一次，\OriginalTerm{物质}{matter}产生了要到达上层区域的欲望；当它到达那里时，它赞叹与天主同在的光辉与光明。的确，它想要为自己夺取优越的位置，并把天主从他的位置挪开。天主还思量怎样报复物质，却缺乏必要的恶来施行，因为恶不存在于天主的居所和住处中。因此，他派遣我们称为灵魂的能力进入物质，让它彻底渗透物质。因为当这能力以后最终与物质分离时，便是物质的死亡。所以，藉着天主的眷顾，灵魂与物质混合在一起，一个不相类者与另一个不相类者相混合。但藉由这混合，灵魂沾染了恶，并陷于与物质同样的病弱。因为正如在败坏的器皿中，所盛之物往往品质变坏，同样，在物质中的灵魂也遭受某种类似的变化，从其本性上恶化，以致有分于物质的恶。但天主怜悯了灵魂，又派遣出另一种能力，我们称之为\OriginalTerm{得穆革}{Demiurge}（即万物的\OriginalTerm{创造者}{Creator}）。当这能力到来，着手创造世界时，它从混合中分离出尽可能多的未曾沾染邪恶和污点的能力，由此首先形成了太阳和月亮；而那沾染了轻微和中等污点的，则成了星辰和天空。从太阳和月亮分离出来的物质，有一部分被完全抛弃在世界之外，就是那火，其中确实有燃烧的能力，虽然它本身是黑暗且无光的，极其像夜晚。但在其余的元素中，无论是动物还是植物，神圣的能力都不均匀地混合其中。因此，世界被造成，其中有太阳和月亮，它们掌管着万物的产生与消亡，通过将神圣的德能从物质中分离出来，并传送给天主。\par

% structure: p0009 source=h2 target=section reason=relative-heading-rank; base=h2

\section{第四章 月亮的盈亏；摩尼教关于此的琐屑之谈；他们对人和\Person{基督}的幻梦；他们愚昧的斋戒制度}

祂诚然安排了此事，为向\OriginalTerm{德穆革}{Demiurge}或创造者提供另一种力量，这力量能吸引到太阳的光辉；这事实甚至可以说对盲人都是显明的。因为月亮在盈满时接收那从物质分离出来的德能，并在其增长期间充满这德能而出。但当它满盈时，在其亏缺中，又将德能返还太阳，太阳则回归于天主。当它这样做完之后，它又再次等待从另一个满月接收灵魂向自身的迁移，并以同样的方式接收后，让灵魂过渡到天主那里。这是它不断的工作，且万古不息。在太阳中可见某种类似人形的形象。而物质又野心勃勃地力图从自身造人，将自己的全部德能混合在一起，以便拥有部分灵魂。但人的形态大大帮助人获得比一切其他动物更大的一份神圣德能。因为人是神圣德能的肖像，而基督则是理智。当祂最终从高处而来时，将这种德能的极大部分遣返于天主。最终被钉十字架，以此方式，祂提供了知识，并使神圣德能被钉在物质中。因此，既然天主的旨意与命令是要物质毁灭，他们便戒绝有生命之物，食用蔬菜及一切无感觉之物。他们也戒绝婚姻、\Person{维纳斯}的仪式及生育子女，以免德能通过种族的延续更深地在物质中扎根；他们也不外出，寻求洁净自己脱离德能与物质混合而沾染的污秽。\par

% structure: p0011 source=h2 target=section reason=relative-heading-rank; base=h2

\section{在天主之下的日月崇拜；在希腊寓言中为摩尼教徒寻求支持；摩尼教徒轻视圣经与信仰的权威}

这些是他们所说所想的主要事项。他们尤其尊崇太阳和月亮，并非视之为神，而是视为可能达到天主的途径。但当神圣德能完全分离之后，他们说外在的火将降下，烧尽自己以及物质所遗留的一切。他们中那些受过更好教育、且不无希腊文学知识的人，从自己的资源中教导我们。例如，从仪式和密仪中：他们说，\Person{巴库斯}从母腹中被切割出来，这标志着神圣德能被\Person{提坦}们分割于物质之中；从诗人关于与\Person{巨人}战斗的寓言中，表明了甚至连那些人也不是不知道物质反抗天主的叛逆。我确实不否认，这些事不足以将那些不加审视而接受话语之人的心灵引入歧途，因为此类言谈所造成的欺骗，已将一些与我一起研究哲学的人吸引到了它那边；至于我当以何种方式接触此事加以审查，我确实迷茫。因为他们的假设并不依循任何正当的方法，以致人可以依照这些方法进行检查；也没有任何论证的原理，使我们能看到由此推演出的结论；他们的只是那些仅可谓为哲思者的罕见发现。这些人，占用旧约与新约圣经，尽管他们断定这些是天主所默感的，却从中引出自己的意见；然后，唯有当他们偶然说出或做出任何与这些不符的事时，才认为被驳倒了。那些按希腊人方式哲学思辨者所拥有的论证原理中的中间命题，在他们那里，就是先知的声音。但在这里，所有这一切都被排除了，既然我之前提到的那些事情是在没有任何论证的情况下提出的，既然有必要以合理的方式给出答复，而不是提出其他更可信且可能更具诱惑力的事情，我的尝试是相当麻烦的，并且因此更为艰巨，因为需要提出性质不同的论证。因为对于那些事先未经证明就被这些人说服的人，更精确的论证将逃过他们的注意，如果在说服之际，他们落入同样的手中。因为他们想象自己是从相似来源出发的。因此，需要大量的、极大的勤奋，且确实需要天主，作我们论证的向导。\par

% structure: p0013 source=h2 target=section reason=relative-heading-rank; base=h2

\section{摩尼教徒的两个原理；其自身被反驳；毕达哥拉斯学派关于第一原理的意见；善与恶相反；胜利在善的一边}

他们设定\OriginalTerm{两个本原}{two principles}，天主与\OriginalTerm{质料}{Matter}。若他（\Person{摩尼}）将生成之物与真实存在者分离开来，这一假设并非十分错误，因为质料既非自我创造，也不接受两种相反的质，既是主动的又是被动的；再者，关于\OriginalTerm{创造因}{creative cause}，也不曾提出其他不可言说的理论。然而，天主为造物并不需要质料，因为万有已实体地存在于祂的理智中，就其可能生成而言。但若如他似乎更倾向的，天主之下的真实存在物的\OriginalTerm{无序运动}{unordered motion}即是质料，那么，首先，他无意中设立了另一个创造因（且是恶的），也未察觉由此而来的结论，即：若必须设定天主与质料，则必须为天主设定另一种质料；从而每一创造因都应有其\OriginalTerm{基质}{subject matter}。因此，他将被证明给出了\OriginalTerm{四个第一本原}{four first principles}，而非两个。这一区分亦令人惊叹。因为若他认为善者即是天主，并愿构想与之对立者，为何不像某些毕达哥拉斯学派那样，将恶置于对立面？实际上，由他们谈论两个本原，善与恶，且二者不断争斗，但善占优势，这尚更可容忍。因为若恶占优势，万物都将毁灭。因此，质料本身既非物体，亦非严格意义上的非形体，亦非单纯某物；而是某种\OriginalTerm{不确定者}{indefinite}，通过添加\OriginalTerm{形式}{form}而得以界定；例如，火是棱锥体，气是八面体，水是二十面体，土是立方体；那么，质料如何是元素的这种无序运动呢？实际上，质料本身并不自立，因为若是运动，它便在于被推动者之中；但质料似乎并非此类存在，而毋宁是\OriginalTerm{第一主体}{first subject}，且是无组织的，其他事物由之产生。既然质料是无序运动，那么它是一直与被推动者相结合，还是曾与被推动者分离？因为，若它曾独自存在，便不会存在；因为没有被推动者则无运动。但若它一直在于被推动者中，那么，又将有两个本原——推动者与被推动者。那么，这两个之中，哪一个将被允许作为与天主并列的\OriginalTerm{第一因}{primary cause}而存在？\par

% structure: p0015 source=h2 target=section reason=relative-heading-rank; base=h2

\section{第七章 为运动辩护，不受不规则之指控；圆周运动；直线运动；论生成与毁灭；论变化与影响感觉的性质}

于此论证附增了一段与此颇不相干的附录。因为你大可以谈论不存在的运动。而且，运动的质料是什么？是圆周运动还是直线运动？抑或通过变化的过程，或通过生成与毁灭的过程而发生？圆周运动确实如此有序且和谐，以至被归诸一切受造物的秩序；而在\OriginalTerm{摩尼教体系}{Manichaean system}中，这运动似乎不配受到非难，因为在其间运行着太阳和月亮，他们声称在诸神中只敬拜这两位。但是，至于直线运动：它一旦抵达其本有位置，便也有其尽头。因为属地之物一触及地面，便完全停止运动。而每一动物和植物，一旦达到其极限，便停止生长。因此，这些事物的止息，较之于他们仿佛为质料所编织的那种无尽死亡，更恰当地说，是质料的死亡。然而，通过生成与毁灭过程而生的运动，则不可能设想为与此假设协调一致，因为据他们说，质料是\OriginalTerm{非受生的}{unbegotten}。但如果他们将变化之运动，即他们所称的那种运动，以及因影响感觉的性质而经受改变的运动归于质料，那么值得思考他们何以得出此说。因为这似乎正是他们所主张的主要论点，因为据他们说，正是通过质料，才导致行为方式改变，灵魂中产生邪恶。因为，在变化中，它将总是从开端开始；继而前行，达至中点，如此便将抵达终点。但当它抵达终点，至少若变化是其本质，就不会停息。而它将再次沿着同一路径返回开端，并同样地由此至终点；且永无休止地如此进行。例如，若\OriginalTerm{\GreekText{α}}{\GreekText{α}}与\OriginalTerm{\GreekText{γ}}{\GreekText{γ}}经历变化，而中介为\OriginalTerm{\GreekText{β}}{\GreekText{β}}，则\GreekText{α}通过改变，将抵达\OriginalTerm{\GreekText{β}}{\GreekText{β}}，并从那里继续至\OriginalTerm{\GreekText{γ}}{\GreekText{γ}}。再从端点\OriginalTerm{\GreekText{γ}}{\GreekText{γ}}返回\OriginalTerm{\GreekText{β}}{\GreekText{β}}，它将在某一时刻抵达\OriginalTerm{\GreekText{α}}{\GreekText{α}}；这一过程持续不断。就像在由黑变白的过程中，中间色为\OriginalTerm{暗褐色}{dun}，端点为白色。反之，由白至暗褐，类似地至黑；变化再次从白色开始，进行同一循环。\par

% structure: p0017 source=h2 target=section reason=relative-heading-rank; base=h2

\section{第八章 质料是邪恶的吗？论天主与质料}

就\OriginalTerm{变化}{alteration}而言，\OriginalTerm{物质}{matter}是\OriginalTerm{恶}{evil}的原因吗？以下证明它并不比\OriginalTerm{善}{good}更恶。假设变化的起点是恶，那么变化便是经由中性之物而从恶到达善。但若变化从善开始，则它同样是通过中性之物而进行。无论\OriginalTerm{运动}{motion}是趋向此一极端或彼一极端，方式都是相同的，这一点已充分说明。所有运动皆与\OriginalTerm{量}{quantity}相关；而\OriginalTerm{质}{quality}则是\OriginalTerm{德性}{virtue}与\OriginalTerm{恶习}{vice}的引导。我们知道，这两者在种类上是不同的。然而，\OriginalTerm{天主}{God}与物质是仅有的\OriginalTerm{本原}{principles}吗，抑或还存在着介于二者之间的某一中项呢？因为若无中项，它们便始终保持互不混合的状态。倘若极端混合了，则必须有某一中介之物来连接它们，这是正确的说法。但若别物存在，该物就必定或是\OriginalTerm{形体}{body}或是\OriginalTerm{非形体的}{incorporeal}，如此一来，第三个外加的本原便出现了。首先，倘若我们设想天主与物质完全是非形体的，那么彼此便不在对方之内，除非像语法知识存在于灵魂中那样；但如此理解天主与物质是荒谬的。然而，若像有些人所言的\OriginalTerm{虚空}{vacuum}那样，虚空由此宇宙所包围；那么另一方则无\OriginalTerm{实体}{substance}，因为虚空的实体乃是虚无。但若二者作为\OriginalTerm{偶性}{accident}，首先这便不可能；因为缺乏实体之物不可能存在于任何地点；实体仿若附性所依托的承载者。若两者均为形体，则二者必同为\OriginalTerm{重}{heavy}、或同为\OriginalTerm{轻}{light}、或同为中间；或一重一轻，或一为中间。若二者皆重，则显然它们在轻质之物与中性之物中应是相同的；或者若它们交替转换，则一方将与另一方全然分离。因为重者有其位置，中间者另有一位置，轻者又有另一位置。一者属上，另一者属下，第三者属中。而在每一\OriginalTerm{球形}{spherical figure}中，下部分即是中央；因为由此处至一切高处，乃至最外层的表面，各方向的距离皆相等，且一切沉重的物体皆从各方被携往该处。因此，每当我听说物质无序地运动——因这属其本性——而来到天主的\OriginalTerm{区域}{region}，或来到\OriginalTerm{光与明亮}{light and brightness}以及此类之处时，我不禁发笑。然若一个为形体，另一个为非形体的，则首先，形体才能单独运动。其次，若二者不相混合，则各自按其本性而彼此分离。但若一个与另一个混和，则它们将成为理智、或灵魂、或\OriginalTerm{偶性}{accident}。因为唯有如此，非形体者才能与形体相混合。\par

% structure: p0019 source=h2 target=section reason=relative-heading-rank; base=h2

\section{第九章 摩尼教徒关于物质向天主运动的荒谬幻想；在摩尼教理解中天主是物质叛乱的原因；物质对光与明亮的渴求是善的；天主的善并不因分施而减少}

然而，物质以何种方式，因何原因而被带到天主的区域呢？因为据他们说，低位和黑暗属于物质的本性；而高位和光则逆其本性。因此，我们便归给它一种\OriginalTerm{超自然的运动}{supernatural motion}；所发生的情形，正如有人将一块石头或一团泥土向上抛掷；依靠抛掷者的\OriginalTerm{力量}{force}，物体被略微升高，当抵达高处区域后，又落回原处。那么，是谁将物质提升至高处区域的呢？物质自身，凭其固有的运动，当然不会如此运动。因此，必须有某种力量施加于它，使其被举升，正如石头和泥块的情形。但他们除了天主以外，没有给物质留下任何别的东西。所以，从他们的论证中，显然会得出如下结论：据他们说，天主以力量与必然性将物质举向了自己。然而，若物质为恶，则其\OriginalTerm{欲求}{desire}全然是恶。对恶的欲求是恶，但对善的欲求则全然是善。既然物质渴求了光与明亮，其欲求便非坏的；正如一个生活在恶习中的人后来欲求德性，这并非坏事。相反，若一人本是善的，却去欲求恶，他并非无罪。这就好比有人说，天主欲求那些依附于物质的恶。因为天主的善物，不应被看作如同巨大的财富、大量的地产和大量的黄金，当一人将它们转让给另一人时，自己所留的便减少了。但如果必须为这些事物在心中构造一个肖像，我想人们会举出\OriginalTerm{智慧}{wisdom}和\OriginalTerm{知识}{sciences}为例。因此，智慧既不减损，知识也不减损，具备这些的人若让他人分享，自己也不会遭受损失；同样，如果我们与他们一同承认物质具有那种欲求，那么认为天主吝于将善的欲求给予物质，便是违背理性的。\par

% structure: p0021 source=h2 target=section reason=relative-heading-rank; base=h2

\section{第十章 关于众神的神话；摩尼教徒的教义与此相似：众神之战的荷马式寓言；根据摩尼教的观点，天主内存在着嫉妒和争胜；这些恶习在任何善人身上都找不到，并且应被视为可耻的}

此外，他们在寓言方面远胜那些神话作者，就是那些要么让\OriginalTerm{科厄卢斯}{Coelus}遭受阉割，要么闲谈\OriginalTerm{萨图恩}{Saturn}被其子设下阴谋以篡夺王位；或是那些说萨图恩吞食其子，并被呈现的石像欺骗而未能达到目的的人。因为他们提出的这些事与那些事有什么不同呢？当他们公开谈论天主与\OriginalTerm{物质}{matter}之间的战争，且不像荷马在\WorkTitle{伊利亚特}中那样以神话的方式讲述这些事——在那里他让\Person{朱庇特}因诸神之间的争斗与战争而欢喜，由此隐约表明世界是由不同等的元素形成的，它们彼此适应，或征服或臣服——。我之所以提出这点，是因为我知道这类人，当他们缺乏证明时，就从各处收集诗歌中的段落，从中寻求对他们自己观点的支持。如果他们稍稍思考一下所遇到的东西，就不会是这样。然而，当一切恶都被从诸神的群体中驱逐时，争竞与嫉妒尤其应当被清除。可是这些人把争竞和嫉妒留给天主，因为他们说天主因物质渴望善便对物质设下计谋。但天主是藉着祂所拥有的哪一样东西来渴望向物质复仇呢？事实上，我认为说天主具有\OriginalTerm{单纯的性体}{simple nature}，比他们所提出的更准确。而且，就像在其他事情上一样，阐述这种想法也不容易。因为它既不能简单地、仅用言语来证明，而需要大量的教导和劳苦。但我们都知道这一点：愤怒与暴怒，以及对物质的报复欲望，是那些被如此激动的人的情感。这种事甚至从来不会发生在一个善人身上，使他受其困扰，更不可能与\OriginalTerm{至善}{Absolute Good}相关连。\par

% structure: p0023 source=h2 target=section reason=relative-heading-rank; base=h2

\section{第十一章 摩尼教徒所传之德能；与同等或较少恶相混合的物质之德能}

因此，我们的讨论又回到了其他事情上。因为，他们宣称天主派遣\OriginalTerm{德能}{virtue}进入物质，值得我们思考这个德能，就善而言，相比于天主，是少于祂呢，还是与祂同等？因为如果少于，原因是什么？因为与天主同在的事物不容与物质有任何相通。但善单独是天主的特征，恶单独是物质的特征。然而，如果它与祂同等，那么是什么原因使祂，如同一位君王，发出号令，而它却非自愿地承担此辛劳？此外，关于物质，应当探究，就恶而言，这些德能是否相似或更少。因为如果它们更少，它们就完全是较少的恶。因此，通过与善的相通，它们才变成这样。因为有两样恶，较少的恶显然通过与善的相通才成为它现在的样子。但他们在物质周围不留任何善。因此，又产生另一个问题。因为如果某种别的德能，就恶而言，胜过占优势的物质，它自己就成了主导的原则。因为那更恶的将在其自己的领域中掌权。\par

% structure: p0025 source=h2 target=section reason=relative-heading-rank; base=h2

\section{第十二章 拒绝以德能注入来消灭恶；因为它不带来恶的减少；弃绝芝诺关于世界将被太阳之火焚烧的观点}

但天主派遣德能进入物质，这一说法毫无证据，且全无可能性。然而，这件事应当有其自身的解释。他们说这样做的理由，确实是为了不再有恶，使一切成为善。德能必须与恶混合，就像角力者紧紧搂抱并制服对手一样，通过征服恶，使它不再存在。但我认为，在天主最初构思存在物时，就消灭物质，这远为更尊贵，也更配得天主的卓越。但我想他们不能允许这样，因为发现有一些恶的东西存在，他们称之为物质。但是，事物并不会为了使人承认某些事物变成了更坏的东西而不再如此。必须对此有所觉察，因为这些当前的事物在某种程度上已经遭受了减少，以便我们对未来有更好的希望。因为这对于\Person{齐乌姆的芝诺}的观点是一个很好的回应，他这样论证世界将被火焚毁：\Quote{凡有可燃之物的事物，在它烧尽一切之前，不会停止燃烧；太阳是火，它不会烧尽它所拥有的东西吗？}由此他推想，宇宙将被火焚毁。但据说有个诙谐的人对他说，但我昨天、去年、很久以前都看到，如今也同样看到，太阳并未受到任何损害；而且有理由认为，这会随着时间逐渐发生，以致我们可以相信，在某个时候，整个宇宙会被焚毁。至于\Person{摩尼}的学说，尽管它毫无证据，我认为可以给出同样的回答，即：事物的现状并无任何减少，相反，起初的时代，当兄弟杀害兄弟之时，至今依然如此；同样的战争，以及更多样化的欲望。现在，有理由认为，这些事物，即使不会完全停止，至少应当减少，如果我们想象它们终将停止的话。但是，既然从他们而来的事情相同，我们对他们未来有什么期待呢？\par

% structure: p0027 source=h2 target=section reason=relative-heading-rank; base=h2

\section{第十三章 在星辰星座中绝无恶；摩尼教观点中生命的一切恶都是虚妄，导致生命的消灭；前文已解释了他们关于灵魂从月亮输送到太阳的幻想}

摩尼将什么称为恶呢？至于太阳和月亮，确实毫无欠缺；但对于诸天和星辰，无论他是否说其中有这类东西，以及是什么东西，我们都理应按次第加以考察。然而，不规则性在他们看来是恶，无序运动亦然；但这些事物永远相同，以同一方式运行；没有人会指责任何行星敢于在黄道带中超逾固定期限而迟延；也不会指责任何恒星，仿佛它不固守在同一位置，不以旋转均匀地环绕世界运转，仿佛每百年后退一步。但对于大地，如果他指责某些地方的崎岖不平，或如果领航员因海上的风暴而恼怒；首先，事实上，按他们所想，这些事都有善在其中。因为倘若大地无物生长，一切动物必迅速死亡。而这结果将急速把大量与物质混合的\OriginalTerm{德能}{virtue}送到天主那里，因此需要许多月亮，去容纳那突然接近的巨大数量。对于海，他们也持同样说法。因为死亡实为意外之幸，好使那消亡之物走上最速通往天主的道路。而大地上的战争、饥荒，以及一切趋于毁灭生命之事，被他们极为尊崇。因为凡是善的原因，都当受尊崇。但这些东西是善的原因，因为伴随它们的毁灭，若它们把从那些死亡者身上分离出来的\OriginalTerm{德能}{virtue}传送给天主的话。\par

% structure: p0029 source=h2 target=section reason=relative-heading-rank; base=h2

\section{第14章 埃及人崇拜的有害动物；人凭借技艺为恶者；法律与纪律纠正贪欲与不义；其中无污点的偶然与必然之事}

而且，看来我们一直不知，埃及人正当地崇拜鳄鱼、狮子和狼，因为这些动物比其他动物更强壮，捕食猎物并将其彻底毁灭；也崇拜鹰和隼，因为它们在空中和地上屠杀较弱的动物。但或许，按照他们的说法，人因此受到特别尊崇，因为人最善于以其精巧的发明和技艺制伏大多数动物。而以免他自己在这善中无份，他便成为他人的食物。因此，再者，在他们看来，那些从微小平凡的种子产生出巨大之物的生成物是荒谬的；而他们认为，更恰当的是天主毁灭这些东西，好使神圣的\OriginalTerm{德能}{virtue}迅速摆脱此生固有的困扰。但对于贪欲、不义及此类事物，\Person{摩尼}会问，我们当说什么呢？诚然，对付这些事，纪律和法律前来救助。纪律，确然，以精心的筹划防止此类事在人中间发生；而法律则对任何从事不义行为时被抓住的人施加惩罚。但，那么，若农夫疏于整治土地，为何应归咎于大地呢？因为按正义所属的天主的主权遭受减损，当它的一些部分出产果实而其他部分则否；或者当风吹起时，由于其他原因，有些人从中得益，而另一些人违愿地承受伤害？他们必定必然地不明了偶然之事与必然之事的性质。否则他们不会将这类事视为怪异。\par

% structure: p0031 source=h2 target=section reason=relative-heading-rank; base=h2

\section{第15章 有知觉之物的贪欲和欲望；魔鬼；有知觉的动物；太阳、月亮和星辰亦如此；柏拉图学说，非基督宗教的}

那么，快乐与欲望从何处来呢？因为这些都是他们谈论并憎恶的主要恶事。物质也\Emph{不显现}为任何别的东西。事实上，这些事仅属于有感觉的动物，而且除了有感觉者之外，无物能感知欲望与快乐，这是明显的。因为植物何有苦乐之感？大地、水、空气又有何感？至于魔鬼，若它们确是有感觉的活物，或许因此乐于所设立的祭献，当这些缺乏时便不悦；但对于天主，绝不能想象此类事。所以那些说\Quote{动物为何受苦乐影响？}的人，应先发怨言：\Quote{这些动物为何有感觉，或为何需要食物？}因为如果动物是不死的，它们就会脱离腐朽与增长；就像太阳、月亮和星辰，尽管它们也有感觉。然而，它们却超乎这些事及此类怨言之上。而人，能够感知、判断，且潜在地智慧——因为他有能力变得如此——当他获得属己之物时，却将它践踏脚下。\par

% structure: p0033 source=h2 target=section reason=relative-heading-rank; base=h2

\section{第16章 因有人智慧，并无妨碍他人亦得智慧；德性须借勤奋与学习而获得；借更健全的哲学，人可被引领向善；基督已为所有人开启了德性的共同研究}

一般而言，值得向这些人询问：是否没有人能够成为善，还是人人都可能有此能力？因为若无人是智慧的，那么\Person{摩尼}自己又算什么呢？我不提他不仅称他人为善，还声称他们能够使别人也成为善。但若有一个个体完全是善的，又是什么阻止了所有人都成为善呢？因为对于一个人是可能的，对于所有人也是可能的。而且通过一个人借以变得有德的同样方法，所有人都可以变得有德，除非他们声称这种\OriginalTerm{德性}{virtue}的较大份额被这类人截获了。因此，再问：首先，屈从于操练有什么必要呢（因为甚至在睡眠中我们都可以变得有德），或者这些人唤醒他们的听众去盼望善又有什么原因呢？因为即便与妓女一同在淤泥中打滚，他们也能获得他们本有的善。但若操练、更好的教导和修德之勤勉能使人变得有德，就让所有人都变得如此吧，而他们那经常重复的说法——\OriginalTerm{质料的无序运动}{unordered motion of matter}——便归于无效了。但对他们而言，更好的说法是：智慧是天主赐予人的工具，为的是通过逐步地引领那因他们被赋予感觉，而出于欲望或快乐所产生的善，逐渐达致善，从而使那些源自他们的荒谬远离他们。因为这样，那些自称为德性教师的人，就会因其目的和生活方式而成为效法的对象，并且便会出现巨大的希望：有朝一日，当所有人都变得智慧，诸恶将会止息。在我看来，\Person{耶稣}所考虑的正是这一点；并且为使农夫、木匠、建造者以及其他工匠不致远离善，祂将他们全体召集在一起，举行了一次公共会议，并通过简单易懂的谈话，既提升了他们对天主的感知，也激发了他们对善的渴望。\par

% structure: p0035 source=h2 target=section reason=relative-heading-rank; base=h2

\section{摩尼教关于质料中德性的观念遭到驳斥；若一种德性被造为非质料的，则其余亦为非质料的；质料性的德性乃是被抛弃的观点}

此外，他们凭什么说天主将\OriginalTerm{神圣德性}{divine virtue}送入\OriginalTerm{质料}{matter}呢？因为，如果它一直都存在，而且既不能以为天主先于它存在，也不能以为质料先于它存在，那么，按照\Person{摩尼}的说法，便又有三个第一原理了。或许，稍后还会出现更多个。但如果它是外来的，是后来才有的东西，又怎能没有质料呢？如果他们使它成为天主的一部分，那么，首先，按照这种想法，他们便断言天主是组合的、有形体的。但这既荒谬、也不可能。如果天主在无质料的情况下制造了它，我奇怪他们（无论是这人自己还是他的门徒）竟没有考虑到：如果（如\OriginalTerm{正统}{orthodox}所说，后续的事物在天主恒存的情况下得以存在）天主按自己的自由意志创造了这德性，那么，为何祂不是所有其他被造物的作者，而这些被造物无需任何先在的质料便被造出来了呢？事实上，这种观点的结论显然是荒谬的；但随之而来的结论，我们接着来看。那么，这德性的本性是否就是要将自己扩展到质料中呢？如果它与它的本性相反，那它又怎能与质料混合呢？但如果这符合它的本性，那么它必定一直且始终与质料同在。但如果这是事实，他们又怎能称质料为恶呢，因为从一开始它就与神圣德性相混合了。还有，那在某个时候与之混合的神圣德性，若又自行分离出去，质料又将如何被毁灭呢？因为，它保守并安全地维持那善的东西，并可能为那些它临在的人带来其他之善，这比给他们带来毁灭或其他某种恶更为合理。\par

% structure: p0037 source=h2 target=section reason=relative-heading-rank; base=h2

\section{摩尼教论消亡与内在；这乃是针对\Person{摩尼}而作的\Emph{\OriginalTerm{以人身为据的论证}{ad hominem}}，因为他自己也在质料之中}

这就是他们所作的聪明断言——就是说，正如我们看到身体在灵魂与之分离时死去，同样，当德性离弃了质料，所剩下的质料也将消散和灭亡。首先，他们确实没有认识到，任何存在物都不能被毁灭成非存在。因为非存在是不存在的。然而，当物体分解并经历变化时，它们便消散了；于是一部分归入土，一部分归入气，一部分归入别的什么。此外，他们不记得自己的教义是：质料乃是无序运动。但是，那自身运动、且运动是其本质而非偶然属性的东西——当德性离去时，那在德性降入之前就已存在的东西，怎会合理地说它就不存在了呢？他们也没有看到区别：一切缺乏灵魂的物体都是不能运动的。因为植物也有生长灵魂。但运动本身，而且是无序运动，他们却将之视为质料的本质。然而，更好的说法是：正如一把音调不准的竖琴，通过加上和谐，一切都趋于协和；同样，那神圣德性，当与那种据他们所说乃是质料的无序运动相结合时，便应该为其本有的无序状态添加某种秩序，并应使之始终适合天主的旨意。因为我要问，\Person{摩尼}本人是怎么变得适合论述这些事的，他究竟是在什么时候宣布了这些事呢？因为他们承认他本人就是质料与那进入其中的德性的混合物。那么，如果是这样，他是在无序运动中讲述了这些事情，这种意见肯定是有缺陷的；或者，他是藉着神圣德性来说这些事的，那么这教条便令人怀疑和不确定了；因为，一方面，就神圣德性那一面说，他分有真理；而在无序运动那一面，他便分有另一部分，并变为虚妄。\par

% structure: p0039 source=h2 target=section reason=relative-heading-rank; base=h2

\section{摩尼教徒的第二德性受制于前述难题，并伴有新的荒谬；德性既是主动的也是被动的，是质料的塑造者，且与它紧密结合；\Person{摩尼}将身体分为三部分}

然而，倘若有人说，\OriginalTerm{天主德能}{divine virtue}已经装饰且仍在装饰\OriginalTerm{质料}{matter}，那么这话就说得聪明得多，也更易于使人们对\Person{摩尼}的学说与言论产生信仰。可是，天主却差遣了另一位德能。就前一位德能所说的话，同样可以用来说这一位；而从他们关于第一位德能的教导所衍生出的种种荒谬，也同样可以拿来用于眼前的情形。然而，另一位（德能），谁又会容忍呢？天主为什么不派遣某一个能成就万事万物的德能呢？既然人的心灵能应对万物如此多才多艺，以至于同一个人竟能通晓几何学、天文学、木工技艺之类，那么天主难道就不能找到那在各方面都足够胜任的独一德能，从而无需第一、第二两位吗？而且，为什么一位德能拥有的力量更像创造者，而另一位却如同承受者和接受者，以便与质料混合呢？因为我在这里还是看不出良好秩序的原因，也看不出与之相反的那种过分的原因。如果它曾是恶的，那它就不曾在天主的居所中。因为既然天主是唯一的善，质料是唯一的恶，我们就必然要说，其他事物都属中间性质，仿佛被安置在中间。可是，当他们说一个原因是创造性的，另一个则与质料混合时，那些具有中间性质的事物，却发现是由另一不同的制造者所造，不是吗？因此，或许这就是较晚近的作者们在那部\WorkTitle{\GreekText{περι}}书中所论及的那个\OriginalTerm{先在原因}{primary antecedent cause}吧。\par

\SourceRecord{Translated by James B.H. Hawkins；Ante-Nicene Fathers；Vol. 6.；Edited by Alexander Roberts, James Donaldson, and A. Cleveland Coxe.；Buffalo, NY: Christian Literature Publishing Co.,；1886.；<http://www.newadvent.org/fathers/0618.htm>.}
